\chapter{Ganglion Excision}

\section*{What Is This Surgery?}
Ganglion excision is a surgical procedure aimed at the removal of ganglion cysts, which are benign, fluid-filled sacs that typically develop near joints or tendon sheaths, most commonly in the hand, wrist, and foot. These cysts arise from the herniation of synovial tissue and are filled with a viscous, gelatinous fluid. While many ganglion cysts are asymptomatic and may resolve spontaneously, surgical intervention is warranted when they cause persistent pain, limit joint function, compress surrounding nerves or vessels, or recur after non-surgical treatments such as aspiration. The surgical goal is to achieve complete excision of the cyst along with its stalk, which connects it to the underlying joint or tendon sheath, thereby minimizing the risk of recurrence and restoring normal function to the affected area.

\section*{Alternative Names}
\begin{itemize}
  \item Ganglion cyst removal
  \item Ganglionectomy
  \item Synovial cyst excision
  \item Tenosynovial cyst excision
  \item Wrist ganglion removal
  \item Dorsal wrist ganglionectomy
  \item Volar wrist ganglionectomy
\end{itemize}

\section*{Relevant Surgical Anatomy}
Ganglion cysts most frequently arise from the dorsal or volar aspects of the wrist, the flexor tendon sheaths of the fingers, or the dorsum of the foot. 

At the dorsal wrist, ganglia typically originate from the scapholunate ligament or the dorsal capsule, often in proximity to the extensor tendon compartments. Key anatomical structures to identify and protect include the superficial radial nerve (especially its dorsal sensory branch), which can be easily injured during dissection, and the extensor tendons. 

Volar wrist ganglia commonly arise from the radiocarpal or scaphotrapeziometacarpal joints and are critically located near the radial artery, which must be carefully identified and preserved. The median nerve and its branches are also in close proximity, particularly with larger volar ganglia.

For finger ganglia (mucous cysts or flexor tendon sheath ganglia), the neurovascular bundles running along the sides of the digits are paramount, as injury can lead to permanent numbness or vascular compromise. Foot ganglia, often on the dorsum, require careful dissection to avoid injury to the superficial peroneal nerve and the extensor tendons. 

Regardless of location, the surgeon must meticulously identify the cyst's origin, which is typically a small pedicle connecting to the joint capsule or tendon sheath, and ensure its complete removal to prevent recurrence.

\begin{itemize}
  \item Superficial radial nerve
  \item Median nerve
  \item Radial artery
  \item Extensor tendons
  \item Flexor tendon sheaths
  \item Neurovascular bundles
\end{itemize}

\section*{Surgical Principle \& Rationale}
The surgical principle of ganglion excision revolves around the complete removal of the cyst wall and its pedicle or stalk, which connects it to the underlying joint capsule or tendon sheath. The rationale for this approach is that the stalk serves as the communication point through which synovial fluid enters and refills the cyst. While aspiration may provide temporary relief, it often leaves this communication intact, leading to high recurrence rates. 

By excising the entire cyst along with a small cuff of the originating joint capsule or tendon sheath, the goal is to eliminate the source of fluid accumulation and seal off the communication, thereby achieving a definitive cure. The technical goals include meticulous dissection to preserve surrounding vital structures such as nerves, arteries, veins, and tendons, minimizing scar formation, and restoring full, pain-free function to the affected joint.

\section*{History \& Surgical Evolution}
Ganglion cysts have been recognized since antiquity, with descriptions attributed to Hippocrates and Galen. Early treatments were often crude, including rupture by blunt force, famously known as "Bible therapy" due to the common practice of striking the cyst with a heavy book. While sometimes effective in decompressing the cyst, this method carried risks of skin damage, infection, and high recurrence. 

The advent of modern surgical principles in the 18th and 19th centuries, coupled with improvements in anesthesia and aseptic technique, allowed for more systematic excision. Early surgical attempts, however, still faced challenges with recurrence, as the importance of excising the stalk was not fully appreciated. Over the 20th century, surgeons like Bunnell and Littler refined the technique, emphasizing meticulous dissection under tourniquet control and magnification to identify and remove the entire cyst and its connection to the joint or tendon sheath. 

More recently, arthroscopic techniques have emerged for dorsal wrist ganglia, offering a minimally invasive alternative with potentially faster recovery and smaller scars, though open excision remains the gold standard for many locations and complex cases.

\section*{Clinical Indications}
Ganglion excision is indicated for symptomatic or problematic ganglion cysts. Specific clinical indications include:

\begin{itemize}
  \item Persistent or increasing localized pain, especially with movement or pressure over the cyst.
  \item Functional impairment, such as reduced range of motion, weakness, or difficulty with grip strength in the hand or wrist.
  \item Compression of adjacent neurovascular structures, leading to symptoms like numbness, tingling, paresthesia, or motor weakness.
  \item Recurrence of the cyst after one or more attempts at aspiration and corticosteroid injection.
  \item Significant cosmetic concern, particularly if the cyst is large or in a highly visible area, and is bothersome to the patient.
  \item Uncertainty in diagnosis, where the swelling could potentially represent a more serious underlying pathology requiring histological confirmation.
  \item Rapid increase in cyst size or changes in consistency that raise suspicion for other conditions.
  \item Mechanical interference with daily activities, footwear, or gloves due to the cyst's size or location.
\end{itemize}

\section*{Clinical Positioning}
Ganglion excision is typically positioned as a secondary or definitive treatment in the overall management algorithm for ganglion cysts. For most asymptomatic or minimally symptomatic ganglia, the initial approach is observation, as many cysts resolve spontaneously. 

If symptoms develop or persist, aspiration of the cyst fluid, often combined with corticosteroid injection, is usually the first-line interventional treatment due to its minimally invasive nature and lower risk profile. However, aspiration carries a significant recurrence rate, often exceeding 50\%. 

Surgical excision is then reserved for cases where conservative measures, including observation and aspiration, have failed to provide lasting relief, or for specific indications such as nerve compression, significant functional impairment, or diagnostic uncertainty. In these scenarios, ganglion excision becomes the primary definitive treatment for the cyst itself, offering the highest success rate in terms of symptom resolution and prevention of recurrence.

\section*{Surgical Technique \& Procedure}
Ganglion excision is typically performed as an outpatient procedure under regional anesthesia (e.g., axillary block for upper extremity) or local anesthesia with sedation, though general anesthesia may be used for patient comfort or complex cases. A tourniquet is usually applied to the limb to provide a bloodless field, enhancing visualization of delicate structures.

\begin{itemize}
  \item \textbf{Anesthesia \& Positioning:} The patient is positioned supine with the affected limb extended on a hand table. After appropriate anesthesia, the limb is prepped and draped in a sterile fashion.
  
  \item \textbf{Incision:} A small, transverse or longitudinal incision is made directly over the cyst, following Langer's lines or skin creases where possible to optimize cosmetic outcome.
  
  \item \textbf{Dissection:} The incision is deepened through the subcutaneous tissue. The cyst is carefully identified and meticulously dissected from the surrounding soft tissues, taking extreme care to protect adjacent nerves, arteries, veins, and tendons. Magnification (loupes) is often used to aid in this critical step.
  
  \item \textbf{Stalk Identification \& Excision:} The cyst is traced to its origin, which is typically a narrow stalk connecting it to the underlying joint capsule or tendon sheath. The entire cyst, along with a small cuff of the joint capsule or tendon sheath at its base, is excised to ensure complete removal of the communication point.
  
  \item \textbf{Inspection \& Hemostasis:} The surgical bed is inspected for any residual cyst tissue or other small ganglia. The tourniquet is deflated, and meticulous hemostasis is achieved using electrocautery.
  
  \item \textbf{Closure:} The wound is irrigated, and the subcutaneous tissues and skin are closed in layers using fine sutures.
  
  \item \textbf{Dressing \& Immobilization:} A sterile dressing is applied, often followed by a bulky soft dressing or a short arm splint, particularly for wrist ganglia, to provide comfort and limit early motion for a few days.
\end{itemize}

\section*{Preoperative Workup \& Preparation}
Preoperative workup for ganglion excision is generally straightforward, given its elective and often minor nature. A thorough medical history and physical examination are essential to assess the patient's overall health and identify any comorbidities. 

Standard preoperative laboratory tests, such as a complete blood count and coagulation profile, may be ordered, especially if the patient is on anticoagulant medications or has a history of bleeding disorders. Imaging studies, most commonly ultrasound, are often utilized to confirm the cystic nature of the mass, delineate its precise anatomical relationships to surrounding neurovascular structures and tendons, and rule out other potential diagnoses. Magnetic resonance imaging (MRI) may be reserved for complex or recurrent cases, or when diagnostic uncertainty persists.

\begin{itemize}
  \item \textbf{Patient Education \& Consent:} Detailed discussion of the procedure, potential risks, benefits, and alternatives, followed by informed consent.
  
  \item \textbf{Medication Review:} Adjustment or temporary cessation of anticoagulant or antiplatelet medications as per surgeon's protocol.
  
  \item \textbf{NPO Status:} Fasting from food and drink for 6-8 hours prior to surgery if sedation or general anesthesia is planned.
  
  \item \textbf{Prophylactic Antibiotics:} Generally not routinely indicated for uncomplicated ganglion excisions, but may be considered for patients with specific risk factors for infection.
  
  \item \textbf{Logistics:} Arrange for transportation home, as driving may be restricted post-anesthesia.
\end{itemize}

\section*{Contraindications \& Precautions}
\textbf{Absolute Contraindications:}

\begin{itemize}
  \item Active infection of the skin overlying the ganglion or in the affected limb, which must be treated prior to elective surgery.
  \item Uncorrected severe coagulopathy or bleeding disorder that cannot be safely managed preoperatively.
  \item Patient refusal or inability to provide informed consent.
  \item A rapidly improving or spontaneously resolving asymptomatic ganglion, as surgery is unnecessary.
\end{itemize}

\textbf{Relative Contraindications \& Precautions:}

\begin{itemize}
  \item Poorly controlled systemic medical conditions (e.g., diabetes, cardiovascular disease) that increase surgical risk; these should be optimized preoperatively.
  \item Proximity of the ganglion to critical neurovascular structures, which necessitates extreme caution and may favor alternative, less invasive treatments if the risk of injury is deemed too high.
  \item Very small, asymptomatic cysts that do not cause functional impairment or cosmetic concern.
  \item A suspected diagnosis other than a simple ganglion (e.g., lipoma, giant cell tumor of tendon sheath, vascular malformation) that requires further diagnostic workup before surgical intervention.
\end{itemize}

Precautions during surgery include the routine use of a tourniquet for a bloodless field, magnification (surgical loupes or microscope), and meticulous, sharp dissection to protect vital surrounding structures.

\section*{Complications \& Risks}
While ganglion excision is generally a safe procedure, potential complications and risks exist, both general surgical and procedure-specific.

\textbf{General Surgical Complications:}

\begin{itemize}
  \item \textbf{Bleeding or Hematoma:} Accumulation of blood under the skin, potentially requiring drainage.
  \item \textbf{Infection:} Wound infection, ranging from superficial cellulitis to deeper abscess formation.
  \item \textbf{Pain:} Postoperative pain, usually mild and manageable with oral analgesics.
  \item \textbf{Scarring:} Formation of a visible scar, which can sometimes be hypertrophic or keloid.
  \item \textbf{Anesthetic Risks:} Adverse reactions to local, regional, or general anesthesia, including nausea, vomiting, allergic reactions, or cardiopulmonary events.
\end{itemize}

\textbf{Procedure-Specific Complications:}

\begin{itemize}
  \item \textbf{Recurrence of the Cyst:} The most common specific complication, occurring in 5-15\% of cases, even after meticulous excision.
  \item \textbf{Nerve Injury:} Damage to adjacent sensory nerves (e.g., superficial radial nerve, digital nerves) leading to numbness, tingling, or neuropathic pain. Motor nerve injury, though less common, can result in weakness.
  \item \textbf{Vascular Injury:} Damage to nearby arteries (e.g., radial artery for volar wrist ganglia) or veins, potentially requiring repair.
  \item \textbf{Tendon Injury or Adhesions:} Damage to tendons or the formation of adhesions between tendons and the surgical site, leading to restricted movement.
  \item \textbf{Joint Stiffness or Reduced Range of Motion:} Particularly after wrist ganglion excision, due to immobilization or scar tissue formation.
  \item \textbf{Complex Regional Pain Syndrome (CRPS):} A rare but severe chronic pain condition that can develop after limb trauma or surgery.
  \item \textbf{Persistent or New Pain:} Pain that does not resolve or new pain developing in the surgical area.
  \item \textbf{Hypersensitive Scar:} The surgical scar may remain tender or sensitive to touch.
  \item \textbf{Incomplete Excision:} Leaving behind a portion of the cyst or its stalk, increasing the risk of recurrence.
\end{itemize}

\section*{Postoperative Recovery Timeline}
Immediately after ganglion excision, the patient is transferred to the Post-Anesthesia Care Unit (PACU) for monitoring of vital signs, pain control, and neurovascular status of the affected limb. The limb is typically elevated to minimize swelling, and ice packs may be applied. Most patients experience mild to moderate pain, which is usually well-controlled with oral analgesics.

During the first 24-48 hours, the patient is discharged home, usually on the same day. The dressing should be kept clean and dry. Gentle movement of unaffected digits is encouraged to prevent stiffness, but strenuous activity and heavy lifting are restricted. Swelling and bruising are common but gradually subside. 

Over the first 1-2 weeks, sutures are typically removed at a follow-up visit, usually between 7 and 14 days post-surgery. Patients are advised to gradually increase activity, avoiding direct pressure on the wound and heavy gripping or lifting. A splint, if applied, is usually worn for a few days to a week to provide support and comfort. Long-term recovery involves continued scar care and a progressive return to full activities, which typically occurs within 4 to 6 weeks, depending on the cyst's size, location, and the individual's healing response.

\section*{Surgical Findings \& Pathology}
During ganglion excision, the surgeon documents the precise location, size, and consistency of the cyst, its origin from the joint capsule or tendon sheath, and its relationship to surrounding neurovascular structures and tendons. The cyst typically appears as a smooth, translucent, thin-walled sac filled with a clear, gelatinous, mucinous fluid. The stalk connecting it to the underlying joint or tendon sheath is carefully identified.

The excised specimen is routinely sent for pathological examination. The pathology report will confirm the diagnosis of a ganglion cyst. Grossly, the specimen is described as a soft, translucent, multiloculated or uniloculated cyst containing viscous, clear fluid. Microscopically, a ganglion cyst is characterized by a fibrous capsule without a true synovial or epithelial lining. It consists of collagen fibers arranged in a disorganized fashion, with myxoid degeneration and mucin-filled spaces. This histological confirmation is crucial to rule out other soft tissue tumors, such as lipomas, giant cell tumors of the tendon sheath, or other rare benign or malignant lesions that can mimic a ganglion cyst clinically.

\section*{Expected Postoperative Course}
A normal and successful postoperative course following ganglion excision is characterized by a gradual resolution of the preoperative symptoms, primarily pain and functional limitations. Patients can expect mild to moderate pain in the immediate postoperative period, which should progressively decrease over several days to a week, managed effectively with over-the-counter or prescribed oral analgesics. 

Swelling and bruising around the surgical site are common initially but will subside within the first two weeks. The wound should heal without complications, and the surgical scar will mature and fade over several months. Full range of motion of the affected joint is typically restored as pain and swelling resolve, and patients can gradually return to their normal daily activities, including work and hobbies, within 4 to 6 weeks, with minimal to no residual discomfort.

\section*{Postoperative Care \& Follow-up}
Postoperative care for ganglion excision focuses on wound healing, pain management, and restoration of function.

\begin{itemize}
  \item \textbf{Wound Care:} Keep the surgical dressing clean and dry. Avoid submerging the wound in water until sutures are removed and the incision is fully healed.
  
  \item \textbf{Elevation \& Ice:} Elevate the affected limb and apply ice packs intermittently for the first 24-48 hours to reduce swelling and pain.
  
  \item \textbf{Pain Management:} Use prescribed or over-the-counter pain relievers as needed.
  
  \item \textbf{Activity Restrictions:} Avoid heavy lifting, strenuous activities, and direct pressure on the surgical site for several weeks as advised by the surgeon. Gentle range-of-motion exercises for unaffected joints are encouraged early.
  
  \item \textbf{Follow-up Appointment:} A follow-up visit is typically scheduled within 1-2 weeks for wound check and suture removal.
  
  \item \textbf{Physical Therapy:} May be recommended in some cases, particularly if significant stiffness or weakness is anticipated, to aid in regaining full range of motion and strength.
  
  \item \textbf{Warning Signs:} Patients should contact their surgeon immediately if they experience increasing pain, redness, swelling, warmth, pus drainage from the wound, fever, new or worsening numbness/tingling, or signs of recurrence.
\end{itemize}

Long-term follow-up primarily involves monitoring for any signs of cyst recurrence, which is the most common specific complication.

\section*{Epidemiology}
Ganglion cysts are the most common soft-tissue tumors of the hand and wrist, accounting for approximately 50-70\% of all soft-tissue masses in this region. They can occur at any age but are most prevalent in individuals between 20 and 40 years old, with a slight female predominance, often cited as a 2:1 or 3:1 female-to-male ratio in many series. 

The most frequent anatomical location is the dorsal aspect of the wrist (60-70\% of cases), followed by the volar aspect of the wrist (15-20\%), and the flexor tendon sheath of the fingers (10-15\%). Ganglia also occur on the dorsum of the foot, knee, and shoulder, though less commonly. Many ganglia are asymptomatic and may resolve spontaneously, meaning the true prevalence in the general population is likely higher than the incidence of cases presenting for medical attention. 

While generally benign, their high incidence contributes significantly to orthopedic and hand surgery caseloads. Mortality associated with ganglion cysts or their excision is negligible, but morbidity can arise from recurrence, nerve injury, or joint stiffness.

\section*{Outcomes \& Prognosis}
The prognosis following surgical excision of a ganglion cyst is generally excellent, with high rates of symptom resolution and patient satisfaction. Success rates, defined as complete relief of pain and restoration of function without recurrence, typically range from 85\% to 95\%. 

Functional outcomes are overwhelmingly positive, with most patients regaining full, comfortable range of motion and strength in the affected limb. Survival outcomes are unaffected, as ganglion cysts are benign lesions. The primary concern and most common reason for reoperation is recurrence of the cyst, which is reported in 5-15\% of cases after open excision, and potentially higher for arthroscopic techniques or incomplete removal of the stalk. 

Prognostic factors influencing recurrence include the completeness of the excision, particularly the removal of the stalk and a cuff of the joint capsule/tendon sheath, and the location of the cyst (some studies suggest volar wrist ganglia may have slightly higher recurrence rates). Overall, ganglion excision is considered a highly effective and safe procedure for symptomatic or recurrent ganglia.

\section*{References}
\begin{enumerate}
  \item MedlinePlus. Ganglion cyst removal. U.S. National Library of Medicine; 2024. Available from: https://medlineplus.gov/ency/article/002996.htm
  
  \item Wolfe SW, Hotchkiss RN, Pederson WC, Kozin SH. Green's Operative Hand Surgery. 7th ed. Elsevier; 2017.
  
  \item American Society for Surgery of the Hand. Ganglion Cysts. ASSH Patient Education; 2024. Available from: https://www.assh.org/handcare/condition/ganglion-cyst
  
  \item Sabiston DC, Townsend CM. Sabiston Textbook of Surgery: The Biological Basis of Modern Surgical Practice. 21st ed. Elsevier; 2021.
\end{enumerate}

\clearpage